\section*{Results and Discussion}

In order to construct an initial configuration which gives reliable 
trends, we need to understand the composition of layers, how far apart
to stack the layers, and how to orient them with respect to each other.
To verify our choices for each parameter, we compare our calculations
to experimental small angle X-ray scattering (SAXS), wide angle X-ray
scattering (WAXS), and ionic conductivity measurements.

To understand the composition of the monomer layers, we ran simulations of 
systems created with 4 - 8 monomers per layer. Both the sandwiched and 
parallel displaced configurations were tested. All systems remained
stable for at least 50 ns. Pair correlation functions calculated between 
aromatic rings along the length of the pores indicate that systems prefer to
stay partitioned into layers (Fig~\ref{fig:zdf}). Pore spacings were 
calculated by measuring the distance between the pore centers. Pore centers
were located by averaging the coordinates of sodium ions in their 
respective pores. Statistics were generated using the bootstrapping 
technique. Systems built with 5 monomers in each layer equilibrate to a pore
spacing that is most consistent with the experimental value of 4.12 nm
derived from SAXS measurements (Table~\ref{table:p2p}). The remainder of
this discussion will focus on the analysis of systems built with 5 
monomers per layer.

The experimental 2D WAXS measurement contains five major reflections which
our model intends to reproduce (Fig ~\ref{fig:WAXSexp}). The first is located
at qz = 1.7 \angstrom \textsuperscript{-1}, corresponding to a real space 
spacing of 3.7 \angstrom. The reflection is attributed to pi-pi stacking
between aromatic rings in the direction perpendicular to the membrane plane, or z-axis. 
For simplicity, this reflection will be referred to as R-pi. A weak intensity
line is located at exactly half the qz value of R-pi (qz = 0.85 \angstrom \textsuperscript{-1}), corresponding to a 
real space periodic spacing of 7.4 \angstrom. This reflection has been 
attributed to 2\textsubscript{1} helical ordering of aromatic rings along the
z axis and will therefore be referred to as R-helix. If positions of the 
aromatic rings can be traced by a helix, then for each turn in the helix, 
there should be two aromatic rings. A third major reflection is marked by 
a low intensity ring located at r = 1.4 \angstrom \textsuperscript{-1}. The
real space separation corresponds to 4.5 \angstrom which is characteristic
of the average spacing between packed alkane chains. This reflection will be
called R-alkanes. Within R-alkanes, are four spots of higher intensity which 
will be called R-spots. All are located $\approx 40$ degrees from the qz axis
in their respective quadrants. In many liquid crystal systems this can be
attributed to the tilt angle of the alkane chains with respect to the xy axis. 
The final major reflections correspond to the spacing and symmetry of the d100 
plane which can be related to the distance between pores. The feature, which 
will be called R-pores, is characterized by dots along qz = 0. The spacing
between dots is indicative of the hexagonal symmetry of the packed pores. It
is easiest to interpret the data by radially integrating the 2D data to get 
a 1D curve which has also been done.

The initial distance between layers can influence the approach towards
equilibrium. Equilibrated systems built according to the 3.7 \angstrom
spacing indicated by R-pi are characterized by a defined, cylindrical and
open pore structure. Aromatic rings are arranged in a helical conformation.
We will refer to this system as Basin A (Fig ~\ref{fig:phaseA}). Simulations
of systems built with layers stacked > 4 \angstrom apart results in a pore
structure characterized by high radial disorder, while still maintaining 
partitioning between hydrophobic and hydrophilic regions. This will be 
called Basin B (Fig~\ref{fig:phaseB}). It is apparent that this LLC membrane
may exist in at least two metastable states. The distinct difference in 
pore structure (Fig~\ref{fig:porestructures}) exhibited by each phase will
likely lead to different transport mechanisms. Understanding which phase 
exists experimentally is necessary in order to appropriately study the system.
% BJC: include more examples of the difference between the two basins
% The arrangement of sodium ions (which are closely bound to carbonyl
% head groups) can be well-approximated by a gaussian distribution 
% Like phase A, phase B can form at 280K. 
% The phase is also present when phase A is heated close to its isotropic transition point
% There are distinct differences in the membrane and pore structures between each state (Fig~\ref{fig:porestructures})

We varied the relative interlayer orientation between sandwiched and 
parallel-displaced based on our knowledge of the stability of these two
pi-stacking modes. Short NVT simulations were run with position restraints
on aromatic ring carbon atoms. Simulated X-ray diffraction patterns 
generated from each configuration establish a difference between the two
stacking modes. In each pattern, R-alkanes is present at a distance of
$\approx 1.5$ \angstrom \textsuperscript{-1}. R-pi is also present in each 
pattern, although it appears to intersect R-alkanes because the spacing
between aromatic rings is similar to the alkane chain packing. The difference 
in aromatic ring spacing between experiment and simulation is likely a result
of the inability of GAFF to properly handle pi interactions. The parallel
displaced configuration contains two lines of high intensity with maximum
intensity occurring where they intersect the alkane chain region due to
constructive interference between the two features. The lines are located
where one would expect to see R-helix. As the simulation is progressed, 
the full line begins to fade and the reflection persists only at the 
intersection with the R-alkanes. Most significantly, the line is no
longer present where R-helix exists and, perhaps coincidentally, the 
intersection with R-alkanes occurs where one would expect to see R-spots.

Full comparison of experimental 2D WAXS with simulated X-ray diffraction
patterns produced from equilibrated MD trajectories shows the most consistency
with the sandwiched conformation in Basin A. Systems were built in both the 
parallel displaced and sandwiched conformations with an intial layer spacing of
3.7 \angstrom. A third system was created by stacking layers in the sandwiched 
conformation 5 \angstrom apart in order to guide it towards Basin B. For ease of
reference, the third system will be referred to as the disordered system. The 
three systems were equilibrated according to our procedure with NPT simulations of 
greater than 400 ns. Simulated diffration patterns were generated using 
portions of the trajectory after equilibration. Equilibration was detected
when the distance between pores and the membrane thickness stopped changing.
Simulated diffraction patterns for all three structures are shown in 
Fig~\ref{fig:XRDsim}. The equilibrated parallel displaced structure exhibits 
R-alkanes and R-pores in the expected locations. Due to low resolution, making
out the individual spots of R-pores is challenging and can be validated upon
further inspection of the 3D structure factor. R-pi is also present, however it
intersects R-alkanes because it is at a lower q than expected. This can be 
explained by our model's inability to appropriately represent aromatic interactions.
The rings prefer to stack $\approx 4.1 \angstrom$ apart as opposed to 3.7 \angstrom.
R-helix and R-spots are both missing. The line which intersected both of these
features in the restrained simulations faded upon equilibration which suggests
that monomers in the parallel-displaced conformation may not give rise to the 
experimental WAXS pattern seen. Simulated XRD of the sandwiched configuration 
contains all features except for R-helix. Most notably, R-spots appears in the
expected location, which suggests that there is something intrinsically to
the sandwiched packing which gives rise to such features. Finally, the 
disordered structure exhibits R-alkanes and R-pores. R-helix are R-spots are not
present. Given the evidence, the sandwiched configuration with layers spaced
3.7 \angstrom apart gives rise to an equilibrated structure closest to experiment. 

The spots that appear in the simulated XRD pattern of the sandwiched
conformation are a result of the way alkane tails pack together. Previously,
the spots in the diffraction pattern had been explained as the product of
tilted alkane chains. We measured the tilt angle of the alkane chains and 
showed that our system equilibrates to an average tilt angle close to zero 
degrees (Fig~\ref{fig:tilt}). To understand the origin of the spots, we 
determined which atoms gave rise to the feature. By removing atoms from 
the trajectory and simulating a diffraction pattern, we were able to isolate
the cause of the spots to the tails. Since the tails are mostly flat, we
plotted the centroids of the tails and measured the angle between each centroid 
and its nearest neighbors with respect to the plane of the membrane. The 
distribution of these angles is consistent with the location of the spots.
We reason that monomer tails stacked closely in the sandwiched conformation
are forced to splay apart and pack in between in each other which creates a
nearly hexagonal array of packed tails. 

        %MRS2: Tables should also have specific theses as well, so think about that.  Merely saying ``here is the information'' isn't as useful without a guide through it.
A summary of all experiments and relevant structural parameters are presented in Table ~\ref{table:p2p}

\begin{figure}
	\centering
	\begin{subfigure}{0.3\textwidth}
		\centering
		\includegraphics[width=\linewidth]{placeholder.png}
		\caption{}\label{fig:orderparameter}
	\end{subfigure}
	\begin{subfigure}{0.3\textwidth}
		\centering
		\includegraphics[width=\linewidth]{placeholder.png}
		\caption{}\label{fig:order}
	\end{subfigure}
	\begin{subfigure}{0.3\textwidth}
		\centering
		\includegraphics[width=\linewidth]{placeholder.png}
		\caption{}\label{fig:disorder}
	\end{subfigure}
	\caption{(a) As temperature is increased, Phase A transitions to 
	Basin B, becoming more disordered. An order parameter of 1 indicates perfect ordering, while 0 indicates disorder. (b) The system is ordered, characteristic of Phase A, initially (c) The system becomes disordered, characteristic of Phase B, as temperature is raised beyond 310K}\label{fig:transition}
	%MRS2: maybe provide a bit more information what data will be here?  Order parametr vs time? One panel with order paramter vs. time, one with configurations from 0 and 1 endpoints?
\end{figure}	

\begin{table}
\centering
\begin{tabular}{|c|c|}
\hline
table & here \\
\hline
\end{tabular}
\caption{The pore spacing of the model increases as number of monomers in each layer increases. The pore spacing of a system starting in the sandwiched configuration is systematically lower than that started in an offset configuration. Systems built with 5 monomers per layer in a parallel displaced configuration results in a pore spacing closest to experimental data}\label{table:p2p} 
\end{table}
	
\begin{figure}
\centering
\begin{subfigure}{.45\textwidth}
	\centering
	\includegraphics[width=\linewidth]{placeholder.png}
	\caption{}\label{fig:phaseA}
\end{subfigure}
\begin{subfigure}{.45\textwidth}
	\includegraphics[width=\linewidth]{placeholder.png}
	\centering
	\caption{}\label{fig:phaseB}
\end{subfigure}
\caption{There are clear differences in pore structure between systems built with layers stacked (a) 3.7 \angstrom apart and (b) 5.0 \angstrom apart}
%MRS2: explain theses?  Also, the exact configurations used to start them are probably not that important; many different initial configurations will probably lead to most of them.  It's the thermodyamics that are important.
%BJC: okay, so just rephrase it, worded in a way that doesn't reference the initial configuration/layer spacing?
%BJC: what do you  mean by 'thermodynamics that are important' -- just distinguishing which state is more favorable under what conditions?
\label{fig:porestructures}
\end{figure}
%BJC: I might eliminate the other paragraph (2 before this one), where I justify parallel-displaced vs. sandwiched, and do it all in this paragraph.	
%MRS2: a little hard to say until the data is in.

\begin{figure}
	\centering
	\begin{subfigure}{.3\textwidth}
		\centering
		\includegraphics[width=\linewidth]{placeholder.png}
		\caption{~\label{fig:xrdA}}
	\end{subfigure}
	\begin{subfigure}{.3\textwidth}
		\centering
		\includegraphics[width=\linewidth]{placeholder.png}
		\caption{~\label{fig:xrdexp}}
	\end{subfigure}
        \begin{subfigure}{.3\textwidth}
                \centering
                \includegraphics[width=\linewidth]{placeholder.png}
                \caption{~\label{fig:xrdB}}
        \end{subfigure}
	\caption{Phase A (a) provides a good match to experiment (b). Phase B (c) is missing reflections present in the experimental pattern} % BJC: any more detail?  %MRS2: yes, give distances, and describe what you mean as ``good match''
%BJC2: what is a clear way to describe the spots in the ring. Diagonal reflections? radial spots?
	\label{fig:xrd}
\end{figure}

The model gives reasonable estimates of ionic conductivity for both phases.
Calculated values of ionic conductivity obtained using the Nernst Einstein
relation and Collective Diffusion model are compared in Table~\ref{table:conductivity}.
The two methods agree with each other within error, although the 
uncertainty obtained using the Collective Diffusion model is much higher.
For this reason we will likely only use the Nernst Einstein relation
in future calculations of this type. It is evident that the ordered 
phase has a higher ionic conductivity than the disordered phase. 
Conductivity is enhanced in Phase A because the pores are more open,
increasing the likelihood of site hopping. In all cases, our calculated values
are higher than experimental values, as expected. The real system, 
although mostly aligned and straight, still has a distribution of 
azimuthal angles, lowering the effective ionic conductivity of the bulk 
membrane. The ordering from isotropic to mostly aligned mesophases 
showed an 85 fold increase in ionic conductivity. We would expect additional
gains in a perfect system.

\begin{table}
\centering
\begin{tabular}{ccc}
\toprule
\multicolumn{3}{c}{Calculated Ionic Conductivity \si{\siemens\per\meter}} \\
\hline
Method & Phase A & Phase B \\
\midrule
Nernst Einstein & \num{8e-5} & \num{8e-5} \\
Collective Diffusion & \num{6e-5} & \num{6e-5} \\
Experiment & \num{1.3e-5} & \num{1.3e-5} \\
\bottomrule
\end{tabular}
\caption{Calculated ionic conductivity using Nernst-Einsten and Collective Diffusion agree within error. Both methods give calculated values of ionic conductivity which are higher, but reasonably close to experimental values~\label{table:conductivity}}
%MRS2: I would say that they are reasonable close to experiment, though higher. Also somewhat contradicted in other places.
	\end{table}
